\documentclass[11pt]{article}
\usepackage[utf8]{inputenc}
\usepackage[T1]{fontenc}
\usepackage[margin=1in]{geometry}
\usepackage{amsmath, amssymb}
\usepackage{booktabs}
\usepackage{array}
\usepackage{hyperref}

\begin{document}

\begin{center}
{\LARGE First Exam — Lessons 0 \& 1 (Version B)}\\[0.75em]
\end{center}

\noindent\textbf{Time allowed:} 60 minutes\\
\textbf{Resources:} Printed course notes and personal handwritten notes only. Electronic devices, calculators, and statistical tables are not permitted.\\
\textbf{Instructions:} Answer all four questions. Show algebraic steps and clearly label each part. Unless stated otherwise, provide exact values (fractions, radicals, exponentials). Partial credit is available when reasoning is clearly explained.

\vspace{0.75em}
\noindent\begin{tabular}{@{}p{2.2cm}p{8.2cm}r@{}}
\toprule
\textbf{Question} & \textbf{Topic} & \textbf{Points}\\
\midrule
1 & Probability structures and independence & 25\\
2 & Bayes' rule and conditional reasoning & 25\\
3 & Uniform model derivations & 25\\
4 & Interpreting summaries and convergence & 25\\
\midrule
\textbf{Total} & & \textbf{100}\\
\bottomrule
\end{tabular}

\vspace{1em}
\hrule
\vspace{1em}

\section*{Question 1 — Probability Foundations (25 points)}

A die is rolled twice. Throughout parts (a)–(c), assume the die is fair.

\begin{enumerate}
  \item (6 pts) Specify a probability space $\,(\Omega, \mathcal{F}, P)\,$ for this experiment. Clearly describe the sample space, a natural $\sigma$-algebra, and the probability measure.
  \item (6 pts) Let $A$ be the event ``the sum of the two rolls is at least 10'' and $B$ the event ``the first roll shows a 5 or 6.'' Compute $P(A)$, $P(B)$, and $P(A \cap B)$.
  \item (5 pts) Are $A$ and $B$ independent? Justify your answer using the formal definition.
  \item (8 pts) Now suppose the die is biased such that the probability of rolling a 6 is $q$ (where $0 < q < \tfrac{1}{6}$), while outcomes 1 through 5 are equally likely with total probability $1-q$. With the same definition of event $B$, derive an expression for $P_q(B)$ in terms of $q$. For which value of $q$ does $P_q(B) = \tfrac{1}{3}$? Provide algebraic reasoning.
\end{enumerate}

\vspace{0.75em}
\hrule
\vspace{0.75em}

\section*{Question 2 — Bayes' Rule in Fraud Detection (25 points)}

A bank uses an automated system to flag potentially fraudulent credit card transactions. Historical data indicate that $\tfrac{1}{50}$ of transactions are actually fraudulent. The system has sensitivity (true positive rate) $\tfrac{4}{5}$ and specificity (true negative rate) $\tfrac{39}{40}$.

\medskip
\noindent\textit{Recall:} sensitivity $= P(\text{Flag} \mid \text{Fraud})$ and specificity $= P(\text{No Flag} \mid \text{No Fraud})$. The false positive rate is $1-\text{specificity} = P(\text{Flag} \mid \text{No Fraud})$.

\begin{enumerate}
  \item (7 pts) Compute the probability that a randomly selected transaction is flagged as suspicious.
  \item (8 pts) Given that a transaction is flagged, determine the probability that it is actually fraudulent. Express your answer as a reduced fraction.
  \item (6 pts) Given that a transaction is not flagged, determine the probability that it is actually fraudulent. Express your answer as a reduced fraction.
  \item (4 pts) Let $\pi = P(\text{Fraud})$, $s = P(\text{Flag} \mid \text{Fraud})$ (sensitivity), and $c = P(\text{No Flag} \mid \text{No Fraud})$ (specificity). Derive a formula for $P(\text{Fraud} \mid \text{Flag})$ in terms of $\pi$, $s$, and $c$, and briefly state how increasing $c$ affects this probability.
\end{enumerate}

\vspace{0.75em}
\hrule
\vspace{0.75em}

\section*{Question 3 — Uniform Waiting-Time Model (25 points)}

Let $T$ be the time (in minutes) until the next customer arrives at a service counter, and suppose $T$ follows a uniform distribution on $[0,8]$.

\medskip
\noindent\textbf{Recall:} The probability density function of a $\mathrm{Uniform}(a,b)$ random variable is
\[
f_T(t) =
\begin{cases}
\dfrac{1}{b-a}, & a \le t \le b,\\
0, & \text{otherwise},
\end{cases}
\]
and the cumulative distribution function is $F_T(t) = 0$ for $t<a$, $F_T(t) = \dfrac{t-a}{b-a}$ for $a \le t \le b$, and $F_T(t) = 1$ for $t>b$.

\begin{enumerate}
  \item (5 pts) Specialise the CDF $F_T(t)$ for $T \sim \mathrm{Uniform}(0,8)$.
  \item (5 pts) Using your result from part (1), compute $P(2 \le T \le 6)$.
  \item (9 pts) Compute $E[T]$ and $\operatorname{Var}(T)$ directly from integrals. Show the essential integration steps leading to exact values.
  \item (6 pts) Define $Y = 2T + 5$. Compute $E[Y]$ and $\operatorname{Var}(Y)$ using properties of expectation and variance. Provide exact values.
\end{enumerate}

\vspace{0.75em}
\hrule
\vspace{0.75em}

\section*{Question 4 — Interpreting Summaries and Convergence (25 points)}

Table~\ref{tab:b-sim} summarises the results of a simulation: for each sample size $n$, 800 independent samples of size $n$ were drawn from an exponential distribution with mean $3$, and the sample mean was recorded for each run.

\begin{table}[h]
\centering
\caption{Behaviour of Sample Means}\label{tab:b-sim}
\begin{tabular}{@{}rcc@{}}
\toprule
\textbf{Sample size $n$} & \textbf{Average of the 800 sample means} & \textbf{SD of the 800 sample means}\\
\midrule
10  & 3.06 & 0.94 \\
40  & 3.01 & 0.47 \\
160 & 3.00 & 0.24 \\
\bottomrule
\end{tabular}
\end{table}

Another analysis compared a standardised dataset (mean 0, variance 1) to the standard normal distribution. Selected points from the QQ-plot are listed below.

\begin{table}[h]
\centering
\caption{Selected QQ-Plot Quantiles}\label{tab:b-qq}
\begin{tabular}{@{}rc@{}}
\toprule
\textbf{Normal quantile} & \textbf{Sample quantile}\\
\midrule
-2.0 & -1.4 \\
-1.0 & -0.8 \\
 0.0 &  0.1 \\
 1.0 &  0.9 \\
 2.0 &  1.6 \\
\bottomrule
\end{tabular}
\end{table}

\medskip
Answer the following:
\begin{enumerate}
  \item (9 pts) Use Table~\ref{tab:b-sim} to explain how the Law of Large Numbers manifests in this simulation. Refer explicitly to the numerical values.
  \item (8 pts) Still relying on Table~\ref{tab:b-sim}, discuss how the variability of the sample mean changes with $n$ and connect your explanation to the Central Limit Theorem.
  \item (8 pts) Interpret Table~\ref{tab:b-qq}. Describe the shape of the corresponding QQ-plot, assess the plausibility of a normal model, and recommend a modelling or diagnostic action before relying on normal-based methods.
\end{enumerate}

\vspace{0.75em}
\hrule
\vspace{0.25em}

\noindent\textit{End of exam.} Ensure answers are clearly labelled and presented legibly.

\end{document}
